\documentclass[11pt,a4paper]{article}

\usepackage[utf8]{inputenc}
\usepackage[T1]{fontenc}
\usepackage{lmodern}
\usepackage{amsmath, amssymb, amsfonts}
\usepackage{bm}
\usepackage{geometry}
\usepackage{graphicx}
\usepackage{hyperref}
\usepackage{authblk}
\usepackage{cite}
\usepackage{physics}
\usepackage{mathtools}

\geometry{margin=2.4cm}

\title{\textbf{Companion XCIX--C \\ Optimal Probes and Benchmark Suite for RDCC}}
\author[1]{Michael Lehmann}
\affil[1]{Independent Researcher}
\date{April 2026}

\begin{document}
\maketitle

\begin{abstract}
This merged Companion unifies the content of Companions XCIX--C into a single 
document. It presents the optimal-probe framework for RDCC inference, including 
criteria for selecting geometric, morphological, and topological observables, 
and provides the benchmark suite used to validate these probes across 
simulations, noise realizations, and tomographic configurations. Together, 
these components form the operational backbone of RDCC data analysis.
\end{abstract}
\section{Optimal Probes for RDCC Inference}

The RDCC framework induces scale-dependent and morphology-dependent signatures 
in the convergence field. Identifying the optimal probes for extracting these 
signatures is essential for efficient and robust inference. This section 
summarizes the criteria for selecting optimal geometric, morphological, and 
topological observables.

\subsection{Definition of an optimal probe}

An observable $O$ is considered an \emph{optimal probe} of the relaxation 
parameter $\alpha$ if it satisfies:
\begin{enumerate}
    \item \textbf{High Fisher information:}
    

\[
    g_O(\alpha) = 
    \mathbb{E}\!\left[
        (\partial_\alpha \log \mathcal{L}_O)^2
    \right]
    \quad \text{is large};
    \]



    \item \textbf{Low degeneracy:}
    the observable responds to $\alpha$ in a direction orthogonal to existing 
    probes;

    \item \textbf{Stability:}
    the response is robust across noise realizations and tomographic schemes;

    \item \textbf{Computational efficiency:}
    the observable can be measured and modeled at scale.
\end{enumerate}

These criteria ensure that optimal probes contribute uniquely and efficiently 
to the inference pipeline.

\subsection{Geometric probes}

Geometric probes include:
\begin{itemize}
    \item convergence power spectra,
    \item tomographic cross-spectra,
    \item angular correlation functions.
\end{itemize}

Their strengths:
\begin{itemize}
    \item well-understood modeling,
    \item high signal-to-noise on linear and quasi-linear scales,
    \item strong constraints on the amplitude and shape of RDCC distortions.
\end{itemize}

Limitations:
\begin{itemize}
    \item reduced sensitivity on nonlinear scales,
    \item partial degeneracy with noise and baryonic effects.
\end{itemize}

\subsection{Morphological probes}

Morphological probes include:
\begin{itemize}
    \item peak counts,
    \item void statistics,
    \item Minkowski functionals.
\end{itemize}

Their strengths:
\begin{itemize}
    \item strong sensitivity to intermediate-scale distortions,
    \item complementary degeneracy structure relative to geometric probes,
    \item robustness to certain noise patterns.
\end{itemize}

Limitations:
\begin{itemize}
    \item increased modeling complexity,
    \item sensitivity to smoothing and filtering choices.
\end{itemize}

\subsection{Topological probes}

Topological probes include:
\begin{itemize}
    \item persistent homology,
    \item Betti curves,
    \item transport-based topology.
\end{itemize}

Their strengths:
\begin{itemize}
    \item highest sensitivity to nonlinear relaxation effects,
    \item strong breaking of degeneracies,
    \item stable information content across noise realizations.
\end{itemize}

Limitations:
\begin{itemize}
    \item computationally expensive,
    \item requires careful treatment of sampling noise.
\end{itemize}

\subsection{Combined optimal-probe strategy}

The optimal RDCC probe combination is a hybrid of:
\begin{itemize}
    \item geometric probes for large-scale structure,
    \item morphological probes for intermediate-scale sensitivity,
    \item topological probes for nonlinear relaxation signatures.
\end{itemize}

This multi-layered structure forms the foundation of the benchmark suite 
presented in the next section.
\section{Benchmark Suite and Validation}

To evaluate the performance of optimal probes and to ensure the robustness of 
RDCC inference, we construct a comprehensive benchmark suite. This suite 
quantifies the sensitivity, stability, and degeneracy structure of each probe 
and validates their performance across simulations, noise realizations, and 
tomographic configurations.

\subsection{Structure of the benchmark suite}

The benchmark suite consists of:
\begin{itemize}
    \item \textbf{Simulation grid:}  
    A grid of RDCC simulations spanning a range of $\alpha$ values, noise 
    levels, and smoothing scales.

    \item \textbf{Probe library:}  
    A collection of geometric, morphological, and topological observables 
    measured consistently across the simulation grid.

    \item \textbf{Response functions:}  
    Numerical derivatives $\partial_\alpha O$ for each observable $O$, 
    estimated using finite differences.

    \item \textbf{Fisher matrices:}  
    Probe-specific and combined Fisher matrices used to quantify information 
    content and degeneracy structure.

    \item \textbf{Null-direction analysis:}  
    Identification of weakly informative directions in observable space.
\end{itemize}

\subsection{Validation criteria}

Each probe is validated according to:
\begin{enumerate}
    \item \textbf{Sensitivity:}  
    The magnitude and stability of $\partial_\alpha O$ across simulations.

    \item \textbf{Consistency:}  
    Agreement between independent noise realizations and tomographic schemes.

    \item \textbf{Degeneracy breaking:}  
    The ability of the probe to break degeneracies present in geometric 
    observables.

    \item \textbf{Cross-probe alignment:}  
    Stability of the probe’s response direction relative to other probes.

    \item \textbf{Computational feasibility:}  
    Runtime and memory requirements for large-scale surveys.
\end{enumerate}

\subsection{Results of the benchmark suite}

The benchmark suite reveals a clear hierarchy:
\begin{itemize}
    \item \textbf{Geometric probes} perform best on large scales and provide 
    stable baseline constraints.

    \item \textbf{Morphological probes} add significant information on 
    intermediate scales and break several geometric degeneracies.

    \item \textbf{Topological probes} provide the strongest constraints on 
    nonlinear relaxation effects and dominate the high-information channels.

    \item \textbf{Combined probes} achieve near-optimal performance, capturing 
    almost all available information about $\alpha$.
\end{itemize}

\subsection{Implications for RDCC analysis}

The benchmark suite provides:
\begin{itemize}
    \item a validated set of optimal probes,
    \item a robust framework for multi-probe inference,
    \item a quantitative understanding of degeneracy structure,
    \item and a foundation for future survey applications.
\end{itemize}

This completes the operational backbone of RDCC inference and prepares the 
ground for the final conclusions of this merged Companion.
\section{Conclusion}

This merged Companion unifies the optimal-probe framework and the benchmark 
suite for RDCC inference. The optimal-probe analysis identifies the geometric, 
morphological, and topological observables that provide the strongest and most 
complementary constraints on the relaxation parameter $\alpha$. The benchmark 
suite validates these probes across simulations, noise realizations, and 
tomographic configurations, revealing a clear hierarchy of information content 
and degeneracy-breaking power.

Together, these components form the operational backbone of RDCC data analysis:
\begin{itemize}
    \item geometric probes provide stable large-scale constraints,
    \item morphological probes enhance intermediate-scale sensitivity,
    \item topological probes dominate nonlinear relaxation signatures,
    \item and combined probes achieve near-optimal performance.
\end{itemize}

This merged document replaces Companions XCIX--C and provides a unified 
reference for probe selection, validation, and multi-probe inference in the 
RDCC framework.

\begin{thebibliography}{9}

\bibitem{RDCC-Flagship}
M.~Lehmann,
\textit{Relaxation-Driven Cyclic Cosmology (RDCC): Flagship Paper},
(2026).

\bibitem{RDCC-Summary}
M.~Lehmann,
\textit{RDCC Summary Paper: Inference Framework and Key Results},
(2026).

\bibitem{RDCC-Observables}
M.~Lehmann,
\textit{RDCC Numerical Fits and Observational Constraints},
(2026).

\bibitem{RDCC-Topology}
M.~Lehmann,
\textit{Topological and Morphological Probes in RDCC},
(2026).

\bibitem{RDCC-Companions-Opt}
M.~Lehmann,
\textit{RDCC Companion Series: Optimal Probes and Benchmark Suite},
Companions XCIX--C (merged), (2026).

\end{thebibliography}

\end{document}
